% =====================================================================
\chapter{Taylor, hessiano y puntos críticos}
\label{tema:4}
% =====================================================================

El gradiente describe la inclinación de una función; la curvatura la describe el hessiano. Este tema introduce las derivadas de segundo orden, demuestra el desarrollo de Taylor de segundo orden en varias variables, lo usa para clasificar puntos críticos y, como aplicación central, determina para qué tasas de aprendizaje converge el descenso de gradiente sobre una función cuadrática. Se sigue~\cite[§5.7 y §5.8]{mml}, \cite[§4.2 y §4.3.1]{gbc} y la clase~9 de~\cite{mit1802}.

\section{El hessiano}

\begin{definicion}[Hessiano]\label{def:hessiano}
Sea $f\colon U\subseteq\R^{n}\to\R$ con derivadas parciales segundas. El \emph{hessiano} de $f$ en $x$ es la matriz $H_f(x)=\nabla^{2}f(x)\in\R^{n\times n}$ de entradas
\[
\bigl(H_f(x)\bigr)_{ij}=\frac{\partial^{2}f}{\partial x_i\,\partial x_j}(x).
\]
Equivalentemente, $H_f$ es el jacobiano del gradiente: $H_f=J_{\nabla f}$.
\end{definicion}

\begin{teorema}[Schwarz]\label{teo:schwarz}
Si $f\in C^{2}(U)$, entonces $\partial^{2}f/\partial x_i\partial x_j=\partial^{2}f/\partial x_j\partial x_i$ en todo punto de $U$. En consecuencia, $H_f(x)$ es simétrica.
\end{teorema}

La demostración se encuentra en los textos de análisis de varias variables y se omite. Al ser simétrico, el hessiano es diagonalizable en una base ortonormal con autovalores reales (teorema espectral, módulo~3), y esos autovalores son las curvaturas de $f$ en las direcciones de los autovectores~\cite[§4.3.1]{gbc}.

\addtocounter{definicion}{1}%

\begin{observacion}[Tamaño del hessiano]
Para $N$ parámetros el hessiano tiene $N^{2}$ entradas. Con $N=7\cdot 10^{9}$ serían $4{,}9\cdot 10^{19}$ números, algo imposible de almacenar. En la práctica solo se calculan productos hessiano-vector $H v$, que cuestan un múltiplo pequeño de un gradiente, o aproximaciones diagonales.
\end{observacion}

\section{Desarrollo de Taylor de segundo orden}

\begin{teorema}[Taylor de segundo orden]\label{teo:taylor}
Sea $f\in C^{2}(U)$ y sean $x\in U$ y $h\in\R^{n}$ tales que el segmento $\{x+th:t\in[0,1]\}$ está contenido en $U$. Entonces:
\begin{enumerate}
  \item (forma de Lagrange) existe $\tau\in(0,1)$ tal que
  \[
  f(x+h)=f(x)+\nabla f(x)\T h+\tfrac12\,h\T H_f(x+\tau h)\,h;
  \]
  \item (forma local) cuando $h\to 0$,
  \[
  f(x+h)=f(x)+\nabla f(x)\T h+\tfrac12\,h\T H_f(x)\,h+o(\norm{h}^{2}).
  \]
\end{enumerate}
\end{teorema}

\addtocounter{definicion}{1}%

\section{Clasificación de puntos críticos}

\begin{definicion}[Punto crítico]
Un punto $x^{*}\in U$ es \emph{crítico} (o estacionario) si $\nabla f(x^{*})=0$. Es un \emph{mínimo local} si existe $r>0$ con $f(x)\ge f(x^{*})$ para $\norm{x-x^{*}}<r$, y un \emph{mínimo local estricto} si la desigualdad es estricta para $x\ne x^{*}$. Un punto crítico que no es mínimo ni máximo local es un \emph{punto de silla}.
\end{definicion}

\begin{teorema}[Condición necesaria de primer orden]\label{teo:fermat}
Si $f$ es diferenciable en $x^{*}$ y $x^{*}$ es un mínimo local, entonces $\nabla f(x^{*})=0$.
\end{teorema}

\begin{teorema}[Criterio del hessiano]\label{teo:criterio}
Sea $f\in C^{2}(U)$ y $x^{*}$ un punto crítico.
\begin{enumerate}
  \item Si $H_f(x^{*})\succ 0$ (todos sus autovalores son positivos), $x^{*}$ es un mínimo local estricto.
  \item Si $H_f(x^{*})\prec 0$ (todos negativos), $x^{*}$ es un máximo local estricto.
  \item Si $H_f(x^{*})$ tiene autovalores de ambos signos, $x^{*}$ es un punto de silla.
\end{enumerate}
\end{teorema}

\begin{table}[h]
\centering\small
\caption{Clasificación de un punto crítico según los autovalores de $H_f(x^{*})$.}\label{tab:criterio}
\begin{tabular}{@{}lll@{}}
\toprule
Autovalores de $H_f(x^{*})$ & Tipo de punto & Ejemplo en $(0,0)$ \\
\midrule
todos $>0$ & mínimo local estricto & $x^{2}+y^{2}$ \\
todos $<0$ & máximo local estricto & $-x^{2}-y^{2}$ \\
de ambos signos & punto de silla & $x^{2}-y^{2}$ \\
alguno nulo, el resto del mismo signo & no concluyente & $x^{4}+y^{2}$, $-x^{4}+y^{2}$ \\
\bottomrule
\end{tabular}
\end{table}

\begin{observacion}[Criterio en dos variables]
Para $n=2$, el producto de los autovalores es $\det H$ y su suma es $\tr H$. Por ello, si $\det H<0$ el punto es de silla, y si $\det H>0$ es un mínimo cuando $f_{xx}>0$ y un máximo cuando $f_{xx}<0$. Es la forma en que se presenta el criterio en~\cite[clase~9]{mit1802}.
\end{observacion}

\addtocounter{definicion}{1}%

\section{Curvatura y descenso de gradiente}

El \cref{teo:taylor} permite analizar el descenso de gradiente más allá de la \cref{prop:descenso}. El caso modelo es una función cuadrática, que coincide con su propio desarrollo de Taylor de segundo orden.

\begin{proposicion}[Descenso de gradiente sobre una cuadrática]\label{prop:dg-cuadratica}
Sea $f(x)=\tfrac12\,x\T Ax$ con $A$ simétrica y autovalores $0<\lambda_1\le\dots\le\lambda_n$. La iteración $x_{k+1}=x_k-\eta\nabla f(x_k)$ converge al mínimo $x^{*}=0$ desde todo punto inicial si y solo si
\[
0<\eta<\frac{2}{\lambda_n}=\frac{2}{\lambda_{\max}(A)}.
\]
\end{proposicion}

\addtocounter{definicion}{1}%

\begin{observacion}[Número de condición]
El factor de contracción en la dirección $i$ es $\abs{1-\eta\lambda_i}$. Con la elección óptima $\eta=2/(\lambda_1+\lambda_n)$, el peor factor vale $(\kappa-1)/(\kappa+1)$, donde $\kappa=\lambda_n/\lambda_1$ es el \emph{número de condición} de $A$. Cuando $\kappa$ es grande, la tasa de aprendizaje está limitada por la dirección de mayor curvatura y el progreso en la de menor curvatura es muy lento. Es el fenómeno de mal condicionamiento descrito en~\cite[§4.2 y §4.3.1]{gbc}, y la motivación de métodos como Adam (módulo~7).
\end{observacion}

\begin{proposicion}[Paso óptimo a lo largo del gradiente]\label{prop:paso-optimo}
Sea $g=\nabla f(x)\ne 0$ y $H=H_f(x)$ con $g\T Hg>0$. La aproximación de segundo orden $q(\eta)=f(x)-\eta\,g\T g+\tfrac12\eta^{2}g\T Hg$ de $f(x-\eta g)$ es mínima en
\[
\eta^{*}=\frac{g\T g}{g\T Hg}.
\]
Si $f$ es cuadrática, $\eta^{*}$ minimiza exactamente $f(x-\eta g)$.
\end{proposicion}

\addtocounter{definicion}{1}%

\begin{observacion}[Relación con la física]
Cerca de un mínimo, $f(x^{*}+h)\approx f(x^{*})+\tfrac12h\T Hh$ es la energía potencial de un sistema de muelles con matriz de rigidez $H$. Los autovalores grandes corresponden a muelles rígidos, que oscilan deprisa; un paso de descenso demasiado largo en esa dirección sobrepasa el equilibrio y el sistema discreto diverge, como establece la \cref{prop:dg-cuadratica}.
\end{observacion}

